\section{Complex-scalar completion and its support graph}
\label{sec:setup}

Let \(I\) and \(J\) be finite sets, let \(b\in\C^I\), and let
\(C=(c_{ij})_{i\in I,j\in J}\in\C^{I\times J}\).  We use
\[
 \langle x,y\rangle=\sum_i\overline{x_i}y_i,
 \qquad
 [x,y]_{\R}=\Real\langle x,y\rangle.
\]
The component value is
\begin{equation}\label{eq:gamma}
 \Gamma(b,C)
 =
 \min_{\theta\in\C^J}\|b+C\theta\|_1.
\end{equation}
Since this is distance to a finite-dimensional subspace, the minimum is
attained.  Standard finite-dimensional convex duality
\cite{Rockafellar1970} gives
\begin{equation}\label{eq:dual}
 \Gamma(b,C)
 =
 \max\left\{
 [z,b]_{\R}:
 z\in\C^I,\ \|z\|_\infty\le1,\ C^*z=0
 \right\}.
\end{equation}
Both extrema are attained.  The condition \(C^*z=0\) is a complex
linear balance equation because the variables in \eqref{eq:gamma}
range over \(\C^J\).

\begin{definition}[Actual incidence graph]\label{def:actual-graph}
The actual bipartite support graph is
\[
 G(C)=(I\sqcup J,E),
 \qquad
 (i,j)\in E\ \Longleftrightarrow\ c_{ij}\ne0.
\]
It is built after all algebraic sums defining \(C\) have been
performed.
\end{definition}

Permutation of rows and columns by the connected components of \(G(C)\)
block diagonalizes \(C\).  Therefore \(\Gamma\) is the sum of its
connected component values, together with \(|b_i|\) for isolated
analysis rows.  Isolated variable columns contribute nothing.  We may
accordingly prove the main statements for a connected component and
sum them afterwards.

\begin{remark}[Why scalar blocks are stated explicitly]\label{rem:scalar}
In TPC-77 an edge could carry a general real-linear coordinate map
\(L:\C\to\C\).  Here an edge carries multiplication by a nonzero
complex number.  This ensures that a leaf variable can realize every
complex correction at its adjacent row and that a matching minor is an
ordinary complex determinant.  \Cref{sec:countermodels} shows that
``nonzero real-linear block'' is insufficient.
\end{remark}

\begin{definition}[Forest component and matching number]
A component is a \emph{forest component} when \(G(C)\) is a tree.
For a forest \(G\), let \(\nu(G)\) denote the largest number of
pairwise vertex-disjoint edges.
\end{definition}

We distinguish two types of leaves throughout:
an analysis leaf belongs to \(I\), while a variable leaf belongs to
\(J\).  Only the second is unconditionally free.

